\section{Avoidance Policies}
\label{sec:planners}

Two local planners are compared. Both consume the same estimated obstacle
stream, the same monocular range model, and the same dead-reckoned pose, and
both emit a body-velocity twist onto the same downstream command path, so the
variable under study is the internal representation of the avoidance decision
and nothing else. Neither planner is offered as a contribution: they exist to
give the predictivity study two behaviorally distinct methods to rank.

\subsection{Obstacle Geometry from a Monocular Detection}

Both planners recover metric geometry from an image-space detection under the
same known-height assumption, so any bias in that model is shared and cannot
favour either method. Bearing follows from the horizontal image coordinate,
$\beta = (c_x - 0.5)\,\mathrm{HFOV}$, and range from the apparent height
under a pinhole model,
\begin{equation}
  \hat{d} \;=\; \frac{H_{\mathrm{ref}}}{2\,h_{\mathrm{bbox}}
                \tan(\mathrm{VFOV}/2)},
  \label{eq:monorange}
\end{equation}
with a class-reference height $H_{\mathrm{ref}}$ and a corresponding
class-reference radius. This model is optimistic when the bounding box is
clipped by the image border, which both planners must absorb through their
safety margins rather than through a privileged measurement.

\subsection{Planner \PC{}: Committed Circumnavigation}

The first planner encodes the avoidance decision as a commitment, which is
the property that makes it robust to perception gaps. On detecting an
obstacle ahead within the engagement range, it selects a side, executes a
lateral displacement using the vehicle's sway authority while continuing
past the obstacle, and returns toward the original route once the obstacle
has been passed. The commitment is the mechanism of interest: a momentary
loss of detection is not evidence that the obstacle is gone, so the maneuver
is not abandoned when the detection stream falls silent, and the side
selection is not revisited mid-maneuver. This is a state-machine policy with
implicit memory, and it is close in shape to published committed-bypass
maneuvers~\cite{zhang2022binocula}; we therefore treat it as a system
component rather than a claim.

\subsection{Planner \PD{}: Holonomic Dynamic Window Baseline}

The classical comparator is a Dynamic Window Approach following the
principle of Fox et al.~\cite{fox1997thedynam} --- sample dynamically
reachable commands, forward-simulate them over a short horizon, discard
inadmissible candidates, and maximize a small normalized objective ---
adapted where the vehicle demands it, in the spirit of Eriksen et al.'s
adaptation to AUV dynamics~\cite{eriksen2016amodifie}. The adaptation is not
claimed as novel; it is a fairness requirement, and we state each departure
from the textbook formulation together with the reason it exists.

\subsubsection{Holonomic search space}
The BlueROV2 Heavy produces lateral thrust natively and the committed planner
uses it heavily, so restricting the baseline to the classical
$(v,\omega)$ differential-drive space would handicap it rather than test it.
The sampled command is therefore $\mathbf{c} = [u, v, r]$ in body surge, sway,
and yaw rate.

\subsubsection{Window and response-aware rollout}
Surge is restricted to what is reachable within one control interval,
$u \in [\max(0, u_0 - a\Delta t),\; \min(u_{\max}, u_0 + a\Delta t)]$ with
$\Delta t$ equal to the replanning period, which is the classical definition;
sway and yaw rate span their full authority because the rollout already
enforces how quickly they are physically attained. That rollout is
response-aware: velocities relax first-order toward the candidate command
during integration, using the same time constants as the navigation model.
Constant-command rollouts, which assume commands are states, over-predicted
clearance by 0.3--0.5\,m under a response lag of order 1\,s and produced
grazing passes; this is a genuine underwater-adaptation requirement rather
than an implementation detail, and it mirrors the reason Eriksen et al.\
introduce second-order dynamics into the window.

\subsubsection{Directional stoppability}
Admissibility retains the classical requirement that the vehicle be able to
stop before contact, but applies it directionally. Writing $c(\cdot)$ for
clearance to the inflated obstacle, a candidate is admissible if its path
clearance stays positive and, \emph{when it is still closing at the end of
the horizon}, its terminal clearance satisfies
\begin{equation}
  c_{\mathrm{end}} \;\geq\; \lVert \mathbf{v} \rVert\, t_{\mathrm{react}}
    \;+\; \frac{\lVert \mathbf{v} \rVert^{2}}{2\,a_{\mathrm{brake}}} .
  \label{eq:stop}
\end{equation}
A candidate whose clearance is increasing has already escaped and is not
required to carry braking distance. The omnidirectional form of this test ---
requiring stopping distance from every candidate, including those moving away
--- makes every lateral escape inadmissible whenever the vehicle is near the
safety boundary, because the path-minimum clearance is dominated by the
shared starting position; the planner then stalls in front of the obstacle.
A separate escape rule applies inside the inflated zone, where the admissible
set is the candidates that increase clearance, so the vehicle cannot be
trapped in a permanent stop within its own margin.

\subsubsection{Mission-paced objective}
The objective keeps Fox's structure --- clearance, progress, speed, heading,
smoothness, all normalized --- with one substantive change. Rewarding raw
forward speed makes the baseline cruise at its maximum surge while the
committed planner holds the mission speed, which is a speed mismatch rather
than a planning difference and which also inflates dead-reckoning drift. The
speed term therefore rewards tracking the nominal mission speed, and the
progress term rewards closing the distance to a receding route goal at
exactly the nominal pace, so that both racing and parking score below
mission-paced motion. This pairing of a dynamic-window search with
line-of-sight guidance at cruise speed is standard marine
practice~\cite{eriksen2016amodifie}. Clearance is scored at the rollout end
state and saturated at a fixed distance, because response-limited rollouts
share a long common prefix that makes path-minimum clearance nearly constant
across candidates.

\subsubsection{Bounded obstacle memory}
A camera obstacle leaves the field of view during the pass while remaining
physically relevant, and a memoryless dynamic window consequently turns back
through the obstacle position once it can no longer see it. The node
therefore maintains a bounded rolling obstacle memory in the odometry frame,
holding only past perception and no ground truth. This is the classical
deployment practice of pairing a dynamic window with a persistent rolling
costmap, and it is an information-parity requirement: the committed planner's
state machine already carries equivalent memory. It is not claimed as novel.

When no candidate is admissible the planner publishes a zero twist and logs a
first-class \texttt{no\_admissible\_candidate} event, which is counted as an
outcome rather than hidden.

\subsection{Input Equivalence}

Fairness was audited against the actual subscriptions and parameters of both
nodes rather than asserted. Both planners receive the same obstacle topic,
compute bearing with the same formula, use identical monocular constants,
read the same dead-reckoned pose, capture the route by the same convention,
apply the same nominal-command staleness rule, and operate under the same
speed limits; neither subscribes to ground truth, and the dynamic window's
planning acceleration bound does not exceed the downstream controller's rate
limit. The dynamic window additionally maintains an internal body-velocity
estimate, but it derives that estimate from its own past commands through the
same first-order model used by the navigation filter, so it consumes no
sensor the committed planner lacks.

The parameters of planner \PC{} are unchanged from the original baseline and
were never tuned on the comparison scenarios. Planner \PD{} is tuned by a
pre-declared grid on development scenarios only, under a lexicographic
objective --- zero collisions first, then success, then clearance above
threshold, then efficiency --- with every tuning run archived. Two pre-freeze
amendments to the declared ranges are recorded with their reasons in the
implementation protocol; both strictly improve the baseline's own safety and
success. The resulting configuration was frozen before the comparison
campaign and is reproduced in Table~\ref{tab:dwaparams}; it was not modified
during the campaign. Two rejected settings are worth recording because they
identify the failure modes the objective must avoid: a dominant clearance
weight parks the vehicle several metres in front of a central obstacle at the
point of maximum clearance, and a weak progress weight parks it likewise ---
both are the classical local minimum, and the mission-paced progress term is
what removes them.

%==============================================================================
% Frozen DWA baseline configuration.  Source: docs/SCIENTIFIC_DECISIONS.md
% D-013 (parameter freeze, 2026-08-11) + docs/DWA_IMPLEMENTATION_AND_PROTOCOL.md.
% Status: READY -- these are the FROZEN values the comparison campaign runs.
% Tuned on P-series development scenarios only; never on F/K campaign scenarios.
%==============================================================================
\begin{table}[t]
\centering
\caption{Frozen configuration of the dynamic-window baseline (planner \PD{}),
decided before the comparison campaign and unchanged during it. Tuning used
development scenarios only, under a lexicographic objective (zero collisions,
then success, then clearance, then efficiency); every tuning run is archived.
Obstacle class constants are applied identically to \emph{both} planners.}
\label{tab:dwaparams}
\footnotesize
\setlength{\tabcolsep}{4pt}
\renewcommand{\arraystretch}{1.15}
\begin{tabular}{@{}p{34mm} p{42mm}@{}}
\toprule
\textbf{Parameter} & \textbf{Frozen value} \\
\midrule
\multicolumn{2}{@{}l@{}}{\emph{Objective weights}} \\
$w_{\mathrm{clearance}}$ & 0.5 \\
$w_{\mathrm{progress}}$ & 1.0 \\
$w_{\mathrm{speed}}$ & 0.3 \\
$w_{\mathrm{route}}$ & 0.5 \\
$w_{\mathrm{smooth}}$ & 0.1 \\
Clearance saturation & 2.0\,m \\
Goal lookahead & 4.0\,m \\
Sway penalty (cruise term) & $0.15/v_{\max}$ \\
\midrule
\multicolumn{2}{@{}l@{}}{\emph{Safety}} \\
Safety margin & 0.80\,m ($+$ vehicle radius 0.40\,m) \\
Stoppability & $v\,t_{\mathrm{react}} + v^{2}/(2a_{\mathrm{brake}})$, \newline
  $t_{\mathrm{react}}=1.2$\,s, $a_{\mathrm{brake}}=0.5$\,m/s$^{2}$, \newline
  closing candidates only, $+$ escape rule \\
\midrule
\multicolumn{2}{@{}l@{}}{\emph{Search and rollout}} \\
Horizon / step & 6.0\,s / 0.2\,s \\
Window & one control interval (0.1\,s); surge slew-limited, sway and yaw full authority \\
Sampling $n_u \times n_v \times n_r$ & $7 \times 7 \times 9$ \\
Velocity limits & $u \leq 0.5$, $|v| \leq 0.3$\,m/s, $|r| \leq 0.3$\,rad/s \\
Response constants $\tau_{u,v,r}$ & 1.2 / 1.15 / 0.3\,s \\
\midrule
\multicolumn{2}{@{}l@{}}{\emph{Obstacle handling}} \\
Memory TTL / merge radius & 30\,s / 1.5\,m \\
Class constants (long range) & $H_{\mathrm{ref}}$ 3.5\,m, radius 1.75\,m \\
Class constants (pool scale) & $H_{\mathrm{ref}}$ 0.5\,m, radius 0.25\,m \\
\bottomrule
\end{tabular}
\end{table}


% TODO_PHASE8_FINAL_RESULT: planner-comparison RESULTS (success, clearance,
% ranking) go in Sec. IX only after the campaign completes and is aggregated.
